\documentclass[10pt,a4paper]{article}
\usepackage[utf8]{inputenc}
\usepackage[T5]{fontenc}
\usepackage[vietnamese,english]{babel}
\usepackage[margin=1.3cm,top=0.75cm,bottom=1.0cm,includehead,includefoot,headheight=18pt,headsep=5pt,footskip=16pt]{geometry}
\usepackage{amsmath,amssymb,array,tabularx,enumitem,fancyhdr,tikz,xcolor,needspace,multicol}
\usetikzlibrary{calc,angles,quotes,arrows.meta,babel}

\definecolor{ink}{RGB}{30,55,70}
\definecolor{accent}{RGB}{18,82,112}

\setlength{\parindent}{0pt}
\setlength{\parskip}{0pt}

\newcommand{\headerQuote}{%
  \ifcase\value{page}\relax
    \or Math gives us hope that every problem has a solution% Page 1
    \or Math + You = Champion% Page 2
    \or Where Frustration Becomes Inspiration% Page 3
    \or Turning Confusion into Clarity% Page 4
    \or Don't be an X, be a Y of mathematics% Page 5
    \or Don't get mad. Learn to add% Page 6
    \or The only way to learn mathematics is to do mathematics% Page 7
    \or Small steps every day lead to big results% Page 8
    \or Every expert was once a beginner% Page 9
    \or Success is the sum of small efforts repeated daily% Page 10
    \or Mathematics is the music of reason% Page 11
    \or Believe you can and you're halfway there% Page 12
    \or Math gives us hope that every problem has a solution% Page 13
    \or Math + You = Champion% Page 14
    \or Where Frustration Becomes Inspiration% Page 15
    \or Turning Confusion into Clarity% Page 16
    \or Don't be an X, be a Y of mathematics% Page 17
    \or Don't get mad. Learn to add% Page 18
    \or The only way to learn mathematics is to do mathematics% Page 19
    \or Small steps every day lead to big results% Page 20
    \or Every expert was once a beginner% Page 21
    \or Success is the sum of small efforts repeated daily% Page 22
    \or Mathematics is the music of reason% Page 23
    \or Believe you can and you're halfway there% Page 24
    \or Math gives us hope that every problem has a solution% Page 25
    \or Math + You = Champion% Page 26
    \or Where Frustration Becomes Inspiration% Page 27
    \or Turning Confusion into Clarity% Page 28
    \or Don't be an X, be a Y of mathematics% Page 29
    \or Don't get mad. Learn to add% Page 30
    \else Math gives us hope that every problem has a solution%
  \fi
}

\pagestyle{fancy}
\fancyhf{}
\lhead{\footnotesize\textit{\color{gray!70!black}\headerQuote}}
\rhead{\small\textbf{\color{ink}Compiled by: Dương Gia Bảo}}
\lfoot{}
\rfoot{\small Page \thepage/3}
\renewcommand{\headrulewidth}{0.4pt}
\renewcommand{\footrulewidth}{0pt}

\fancypagestyle{dapan}{
  \fancyhf{}
  \lhead{\footnotesize\textit{\color{gray!70!black}\headerQuote}}
  \rhead{\small\textbf{\color{ink}Compiled by: Dương Gia Bảo}}
  \lfoot{\small\textit{\color{gray!60!black}Answer Key — Grade 10 Math Review (Chapter IV: Vectors — Test 3)}}
  \rfoot{\small Page \thepage/2}
  \renewcommand{\headrulewidth}{0.4pt}
  \renewcommand{\footrulewidth}{0pt}
}

\newcommand{\parttitle}[1]{\par\vspace{3.5pt}\needspace{3\baselineskip}{\large\bfseries\color{ink}#1}\par\nopagebreak\vspace{2pt}}
\newcommand{\four}[4]{\par\noindent\begin{tabularx}{\linewidth}{@{}*{4}{>{\raggedright\arraybackslash}X}@{}}\textbf{A.} #1&\textbf{B.} #2&\textbf{C.} #3&\textbf{D.} #4\end{tabularx}}
\newcommand{\two}[4]{\par\noindent\begin{tabularx}{\linewidth}{@{}*{2}{>{\raggedright\arraybackslash}X}@{}}\textbf{A.} #1&\textbf{B.} #2\\[3.8pt]\textbf{C.} #3&\textbf{D.} #4\end{tabularx}}
\newcommand{\one}[4]{\par\noindent\begin{tabularx}{\linewidth}{@{}>{\raggedright\arraybackslash}X@{}}\textbf{A.} #1\\[4pt]\textbf{B.} #2\\[4pt]\textbf{C.} #3\\[4pt]\textbf{D.} #4\end{tabularx}}
\newcommand{\answerbox}{\hfill\mbox{\textit{Answer:}\ \fbox{\rule{0pt}{7.5pt}\hspace{1.8cm}}}}

\setlist[enumerate,1]{label=\textbf{Question \arabic*.},leftmargin=2.1cm,labelsep=.15cm,itemsep=1.2pt,topsep=1.2pt,parsep=0pt}
\setlist[enumerate,2]{label=\alph*),leftmargin=.55cm,itemsep=1.5pt,topsep=0.2pt,parsep=0pt}
\newcolumntype{C}[1]{>{\centering\arraybackslash}m{#1}}


\begin{document}

\noindent
\begin{minipage}[c]{0.38\textwidth}
    \centering
    \textbf{DO MATH TO BE HAPPY}\\[-2pt]
    \rule{3.2cm}{0.6pt}\\[2pt]
    \textbf{OFFICIAL REVIEW TEST}
\end{minipage}\hfill
\begin{minipage}[c]{0.52\textwidth}
    \centering
    \textbf{GRADE 10 MATH REVIEW — CHAPTER IV. VECTORS (TEST 3)}\\[1pt]
    \textbf{ACADEMIC YEAR 2026 -- 2027}\\[1pt]
    \textbf{Subject: MATHEMATICS -- Grade 10}\\
    \textit{\small Time allowed: 90 minutes (excluding distribution time)}
\end{minipage}

\vspace{0.05cm}
\noindent
\begin{minipage}[b]{0.78\textwidth}
    \textbf{Full Name:}\dotfill\textbf{ID:}\makebox[3.5cm]{\dotfill}
\end{minipage}\hfill
\begin{minipage}[b]{0.20\textwidth}
    \hfill
    \setlength{\fboxrule}{0.8pt}%
    \fbox{\begin{minipage}{2.3cm}
        \centering
        \textbf{\small Test Code}\\
        \textbf{1044-EN}
    \end{minipage}}
\end{minipage}
\vspace{0.05cm}

\parttitle{PART I. Multiple-Choice Questions}
\textit{Answer Questions 1 to 20. Choose only one correct option for each question.}
\begin{enumerate}
\item Given a parallelogram $ABCD$ with center $O$. How many nonzero vectors are collinear (parallel) with $\vec{AB}$ and have both initial and terminal points at the vertices of parallelogram $ABCD$?
\four{$2$.}{$6$.}{$8$.}{$4$.}

\item Given a regular hexagon $ABCDEF$ with center $O$. Which of the following vectors is equal to $\vec{AB}$?
\four{$\vec{CF}$.}{$\vec{OC}$.}{$\vec{DE}$.}{$\vec{OF}$.}

\item For any four points $M, N, P, Q$, which of the following vector equations is always true?
\two{$\vec{MN} + \vec{PQ} = \vec{MQ} + \vec{NP}$.}{$\vec{MN} - \vec{MP} = \vec{NP}$.}{$\vec{MN} + \vec{PQ} = \vec{MQ} + \vec{PN}$.}{$\vec{MN} + \vec{NQ} = \vec{QM}$.}

\item Let $ABC$ be a right triangle at $A$ with $AB = 3$ and $AC = 4$. What is the magnitude of the vector $\vec{u} = \vec{AB} - \vec{AC}$?
\four{$5$.}{$7$.}{$1$.}{$\sqrt{7}$.}

\item Let $ABCD$ be a rhombus with side length $a$ and $\widehat{ABC} = 60^\circ$. What is the magnitude of the vector $\vec{v} = \vec{BA} + \vec{BC}$?
\four{$a$.}{$a\sqrt{3}$.}{$a\sqrt{2}$.}{$2a$.}

\item Let $G$ be the centroid of triangle $ABC$ and $O$ be any point in the plane. Which of the following statements is true?
\two{$\vec{OA} + \vec{OB} + \vec{OC} = \vec{OG}$.}{$\vec{OA} + \vec{OB} + \vec{OC} = 2\vec{OG}$.}{$\vec{GB} + \vec{GC} = \vec{GA}$.}{$\vec{OA} + \vec{OB} + \vec{OC} = 3\vec{OG}$.}

\item Given line segment $AB$ and a point $M$ lying on segment $AB$ such that $MA = 3MB$. Which of the following vector equations is correct?
\two{$\vec{AM} = \dfrac{3}{4}\vec{AB}$.}{$\vec{AM} = \dfrac{1}{4}\vec{AB}$.}{$\vec{AM} = -\dfrac{3}{4}\vec{AB}$.}{$\vec{AM} = \dfrac{4}{3}\vec{AB}$.}

\item Given a parallelogram $ABCD$, let $M$ be the midpoint of side $BC$. Expressing vector $\vec{AM}$ in terms of $\vec{AB}$ and $\vec{AD}$ gives
\two{$\vec{AM} = \dfrac{1}{2}\vec{AB} + \vec{AD}$.}{$\vec{AM} = \vec{AB} - \dfrac{1}{2}\vec{AD}$.}{$\vec{AM} = \vec{AB} + \dfrac{1}{2}\vec{AD}$.}{$\vec{AM} = \dfrac{1}{2}\vec{AB} + \dfrac{1}{2}\vec{AD}$.}

\item In the Cartesian coordinate plane $Oxy$, given point $A(2, -1)$ and vector $\vec{u} = -3\vec{i} + 4\vec{j}$. What are the coordinates of vector $\vec{v} = \vec{OA} + \vec{u}$?
\four{$(5, -5)$.}{$(1, -3)$.}{$(-1, -5)$.}{$(-1, 3)$.}

\item In the coordinate plane $Oxy$, given two vectors $\vec{a} = (1, -2)$ and $\vec{b} = (-3, 5)$. What are the coordinates of vector $\vec{c} = 2\vec{a} - \vec{b}$?
\four{$(5, -9)$.}{$(-1, 1)$.}{$(5, 1)$.}{$(-5, 9)$.}

\item In the coordinate plane $Oxy$, given triangle $ABC$ with $A(1, 4)$, $B(-3, 2)$, and $C(5, -3)$. What are the coordinates of the centroid $G$ of triangle $ABC$?
\four{$G(3, 3)$.}{$G(1, 1)$.}{$G(1, 3)$.}{$G(-1, 1)$.}

\item In the coordinate plane $Oxy$, given three points $A(-1, 2)$, $B(2, -4)$, and $C(m, 4)$. For what value of parameter $m$ are points $A, B, C$ collinear?
\four{$m = 2$.}{$m = 0$.}{$m = -2$.}{$m = -4$.}
\pagebreak

\item In the coordinate plane $Oxy$, given two vectors $\vec{u} = (m, -4)$ and $\vec{v} = (-3, 6)$. For what value of $m$ are vectors $\vec{u}$ and $\vec{v}$ parallel (collinear)?
\four{$m = 2$.}{$m = -2$.}{$m = 8$.}{$m = -8$.}

\item Let $ABC$ be an isosceles right triangle at $A$. What is the angle between the vectors $\vec{AB}$ and $\vec{BC}$?
\four{$45^\circ$.}{$90^\circ$.}{$60^\circ$.}{$135^\circ$.}

\item Let $ABC$ be an equilateral triangle with side length $a$. What is the dot product $\vec{AB} \cdot \vec{AC}$?
\four{$a^2$.}{$\dfrac{a^2}{2}$.}{$-\dfrac{a^2}{2}$.}{$\dfrac{a^2\sqrt{3}}{2}$.}

\item In the coordinate plane $Oxy$, given two vectors $\vec{a} = (2, -1)$ and $\vec{b} = (3, 4)$. What is the value of $\vec{a} \cdot (\vec{a} + \vec{b})$?
\four{$2$.}{$5$.}{$7$.}{$13$.}

\item In the coordinate plane $Oxy$, given two points $A(1, -2)$ and $B(-5, 6)$. What is the length of segment $AB$?
\four{$10$.}{$100$.}{$2\sqrt{13}$.}{$14$.}

\item In the coordinate plane $Oxy$, given two vectors $\vec{u} = (m + 1, 3)$ and $\vec{v} = (2, -4)$. For what value of $m$ is $\vec{u} \perp \vec{v}$?
\four{$m = -5$.}{$m = 7$.}{$m = -7$.}{$m = 5$.}

\item In the coordinate plane $Oxy$, given two vectors $\vec{x} = (1, \sqrt{3})$ and $\vec{y} = (\sqrt{3}, 1)$. What is the angle between $\vec{x}$ and $\vec{y}$?
\four{$60^\circ$.}{$30^\circ$.}{$45^\circ$.}{$120^\circ$.}

\item A person pushes a cart over a distance of $20\text{ m}$ using a constant force $\vec{F}$ of magnitude $80\text{ N}$ making an angle of $30^\circ$ with the direction of motion. What is the work $W$ done by force $\vec{F}$?
\four{$1600\text{ J}$.}{$800\text{ J}$.}{$800\sqrt{3}\text{ J}$.}{$1600\sqrt{3}\text{ J}$.}
\end{enumerate}

\vspace{3pt}
\parttitle{PART II. True/False Questions}
\textit{Answer Questions 1 to 4. For each statement a), b), c), d) in every question, indicate whether it is True or False.}
\begin{enumerate}
\item Given an equilateral triangle $ABC$ with side length $a$, $M$ is the midpoint of side $BC$, and $G$ is the centroid of triangle $ABC$ (as shown in the figure).
\par\noindent\begin{minipage}[t]{0.72\linewidth}
\begin{enumerate}
    \item The magnitude of vector $\vec{AG}$ is $\dfrac{a\sqrt{3}}{3}$.
    \item The vector equality $\vec{AB} + \vec{AC} = 3\vec{AG}$ is true.
    \item The dot product is $\vec{GA} \cdot \vec{GB} = \dfrac{a^2}{6}$.
    \item The magnitude of vector $\vec{GB} - \vec{GC}$ is $a$.
\end{enumerate}
\end{minipage}\hfill
\begin{minipage}[t]{0.26\linewidth}\raggedleft\vspace{-6pt}
\begin{tikzpicture}[scale=0.65, >=Stealth]
  \coordinate (B) at (0,0);
  \coordinate (C) at (3,0);
  \coordinate (A) at (1.5, 2.598);
  \coordinate (M) at (1.5, 0);
  \coordinate (G) at (1.5, 0.866);
  \draw[thick, ink] (A) -- (B) -- (C) -- cycle;
  \draw[dashed, gray!70] (A) -- (M);
  \draw[dashed, gray!70] (B) -- (G) -- (C);
  \fill[ink] (G) circle (1.8pt);
  \node[above, font=\scriptsize] at (A) {$A$};
  \node[below left, font=\scriptsize] at (B) {$B$};
  \node[below right, font=\scriptsize] at (C) {$C$};
  \node[below, font=\tiny] at (M) {$M$};
  \node[above right, font=\scriptsize] at (G) {$G$};
\end{tikzpicture}
\end{minipage}

\item Given triangle $ABC$. Let $M$ be the midpoint of side $BC$ and $N$ be a point on segment $AC$ such that $3NA = 2NC$ (ratio $NA : NC = 2 : 3$).
\begin{enumerate}
    \item Vector $\vec{AM}$ can be expressed as $\vec{AM} = \dfrac{1}{2}\vec{AB} - \dfrac{1}{2}\vec{AC}$.
    \item Vector $\vec{AN}$ can be expressed as $\vec{AN} = \dfrac{2}{3}\vec{AC}$.
    \item Vector $\vec{BN}$ can be expressed as $\vec{BN} = -\vec{AB} + \dfrac{2}{5}\vec{AC}$.
    \item Vector $\vec{MN}$ can be expressed as $\vec{MN} = -\dfrac{1}{2}\vec{AB} - \dfrac{1}{10}\vec{AC}$.
\end{enumerate}

\item In the coordinate plane $Oxy$, given three points $A(1, 2)$, $B(4, 6)$, and $C(9, 6)$.
\begin{enumerate}
    \item Point $D$, symmetric to $A$ across $B$, has coordinates $D(7, 10)$.
    \item Triangle $ABC$ is an isosceles right triangle at $B$.
    \item The length of segment $AC$ is $4\sqrt{5}$.
    \item The dot product is $\vec{AB} \cdot \vec{AC} = 20$.
\end{enumerate}

\item Given rectangle $ABCD$ with $AB = 6$ and $AD = 8$.
\begin{enumerate}
    \item The dot product is $\vec{AC} \cdot \vec{BD} = 0$.
    \item The dot product is $\vec{AB} \cdot \vec{AC} = 36$.
    \item The magnitude of vector $\vec{AC} + \vec{BD}$ is $16$.
    \item The magnitude of vector $\vec{AB} + \vec{AC}$ is $14$.
\end{enumerate}
\end{enumerate}

\pagebreak
\parttitle{PART III. Short-Answer Questions}
\textit{Answer Questions 1 to 6. Write your final answer in the corresponding box.}
\begin{enumerate}
\item Given triangle $ABC$ with $AB = 5$, $AC = 8$, and $\widehat{BAC} = 60^\circ$. Find the magnitude of vector $\vec{u} = \vec{AB} - \vec{AC}$.\answerbox

\item Given triangle $ABC$. Let $M$ be the midpoint of side $BC$ and $I$ be the midpoint of segment $AM$. Given that $\vec{BI} = x\vec{AB} + y\vec{AC}$ where $x, y \in \mathbb{R}$, calculate $P = 8y - 4x$.\answerbox

\item In the coordinate plane $Oxy$, given parallelogram $ABCD$ with $A(-1, 3)$, $B(2, 5)$, and $C(4, 1)$. Let $D(x_D, y_D)$ be the fourth vertex; calculate $T = 4x_D - 5y_D$.\answerbox

\item In the coordinate plane $Oxy$, given three points $A(1, 2)$, $B(4, 6)$, and $C(-2, 6)$. Find the value of $\cos \widehat{BAC}$ (express your answer as an exact decimal).\answerbox

\item In the coordinate plane $Oxy$, given vectors $\vec{a} = (2, 2)$, $\vec{b} = (-1, 3)$, and $\vec{c} = (5, 1)$. Find the value of the real parameter $k$ such that $\vec{u} = \vec{a} + k\vec{b}$ is orthogonal to $\vec{c}$.\answerbox

\item \begin{minipage}[t]{0.70\linewidth}
A worker pulls a crate across a horizontal floor with a rope making an angle of $30^\circ$ above the horizontal (as shown in the figure). Given a constant pulling force of magnitude $F = 40\sqrt{3}\text{ N}$ and a straight-line displacement of $s = 15\text{ m}$, calculate the mechanical work $W$ done by force $\vec{F}$ (in Joules).\answerbox
\end{minipage}\hfill
\begin{minipage}[t]{0.28\linewidth}\raggedleft\vspace{-6pt}
\begin{tikzpicture}[scale=0.55, >=Stealth]
  \draw[thick, gray!80!black] (-0.3,0) -- (4.2,0);
  \foreach \x in {0.0,0.4,...,4.0} {
    \draw[gray!60] (\x,0) -- (\x-0.18,-0.18);
  }
  \fill[accent!15] (0.5,0) rectangle (2.1,1.0);
  \draw[thick, ink] (0.5,0) rectangle (2.1,1.0);
  \node[font=\scriptsize, ink] at (1.3, 0.5) {$m$};
  \coordinate (P) at (2.1, 0.65);
  \coordinate (H) at (4.1, 0.65);
  \coordinate (Fend) at (3.85, 1.66);
  \draw[dashed, gray!70] (P) -- (H);
  \draw[->, very thick, accent] (P) -- (Fend) node[above right, font=\scriptsize] {$\vec{F}$};
  \draw[ink!80] ($(P)+(0.75,0)$) arc (0:30:0.75);
  \node[font=\tiny] at ($(P)+(1.32,0.22)$) {$30^\circ$};
  \draw[->, thick, ink] (0.8, 1.35) -- (2.3, 1.35) node[midway, above, font=\tiny] {$\vec{s}$};
\end{tikzpicture}
\end{minipage}
\end{enumerate}

\vspace{3pt}
\parttitle{PART IV. Free-Response Problems}
\textit{Show complete reasoning and step-by-step solutions for the following problems:}

\vspace{4pt}
\textbf{Problem 1 (1.0 point).} Given triangle $ABC$. Let $I$ be the midpoint of side $BC$, $K$ be the midpoint of segment $AI$, and $D$ be a point on side $AC$ such that $AD = \dfrac{1}{3}AC$.
\begin{enumerate}[label=\alph*), leftmargin=0.6cm, itemsep=2pt]
    \item (0.5 point) Express vectors $\vec{BK}$ and $\vec{BD}$ in terms of $\vec{AB}$ and $\vec{AC}$.
    \item (0.5 point) Using the results from part a), prove that points $B, K, D$ are collinear.
\end{enumerate}

\vspace{6pt}
\textbf{Problem 2 (1.0 point).} In the coordinate plane $Oxy$, given triangle $ABC$ with vertices $A(-2, 6)$, $B(6, 2)$, and $C(-3, -1)$.
\begin{enumerate}[label=\alph*), leftmargin=0.6cm, itemsep=2pt]
    \item (0.5 point) Find the coordinates of point $M$ on the $x$-axis ($Ox$) that minimizes $|\vec{MA} + \vec{MB}|$, and determine that minimum value.
    \item (0.5 point) Find the coordinates of the circumcenter $I$ of triangle $ABC$.
\end{enumerate}

\vspace{6pt}
\textbf{Problem 3 (1.0 point).} Two tugboats simultaneously pull a large ship at point $O$ with pulling forces $\vec{F}_1$ and $\vec{F}_2$ of magnitudes $F_1 = 5000\text{ N}$ and $F_2 = 8000\text{ N}$, respectively, making an angle of $120^\circ$ with each other (as shown in the figure).
\par\noindent\begin{minipage}[t]{0.70\linewidth}
\begin{enumerate}[label=\alph*), leftmargin=0.6cm, itemsep=2pt]
    \item (0.5 point) Express the resultant force vector $\vec{F}$ in terms of $\vec{F}_1, \vec{F}_2$ and compute its magnitude $|\vec{F}|$ acting on the ship.
    \item (0.5 point) Choose a coordinate system $Oxy$ such that the positive $x$-axis ($Ox$) aligns with $\vec{F}_1$ (as shown). Find the coordinates of $\vec{F}_1, \vec{F}_2, \vec{F}$ and determine the angle $\alpha$ between $\vec{F}$ and $\vec{F}_1$ (round to one decimal place in degrees).
\end{enumerate}
\end{minipage}\hfill
\begin{minipage}[t]{0.28\linewidth}\raggedleft\vspace{-15pt}
\begin{tikzpicture}[scale=0.58, >=Stealth]
  \coordinate (O) at (0,0);
  \coordinate (F1) at (1.75, 0);
  \coordinate (F2) at (-1.4, 2.425);
  \coordinate (F) at (0.35, 2.425);
  \draw[->, thin, gray!80!black] (-2.0, 0) -- (2.5, 0) node[right, font=\tiny] {$x$};
  \draw[->, thin, gray!80!black] (0, -0.5) -- (0, 2.9) node[above, font=\tiny] {$y$};
  \fill[ink!85] (-0.35,-0.22) -- (0.25,-0.22) -- (0.45,0) -- (0.25,0.22) -- (-0.35,0.22) -- cycle;
  \draw[->, very thick, accent] (O) -- (F1) node[below, font=\scriptsize] {$\vec{F}_1$};
  \draw[->, very thick, accent] (O) -- (F2) node[above left, font=\scriptsize] {$\vec{F}_2$};
  \draw[dashed, gray!60] (F1) -- (F) -- (F2);
  \draw[->, line width=1.1pt, ink] (O) -- (F) node[above right, font=\scriptsize] {$\vec{F}$};
  \draw[ink!80] (0.42,0) arc (0:120:0.42);
  \node[font=\tiny] at (0.95, 0.34) {$120^\circ$};
  \node[below left, font=\tiny] at (O) {$O$};
\end{tikzpicture}
\end{minipage}


\newpage
\pagestyle{dapan}
\setcounter{page}{1}

\begin{center}
  {\Large\bfseries\color{ink} COMPREHENSIVE ANSWER KEY}\\[2pt]
  {\small\color{gray!75!black} Grade 10 Math Review — Chapter IV. Vectors (Test 3 — Code 1044-EN)}
\end{center}
\vspace{1pt}

\noindent\textbf{PART I. Multiple-Choice Questions (20 questions)}
\vspace{1.5pt}

\begin{center}
\renewcommand{\arraystretch}{1.15}
\begin{tabular}{|C{0.95cm}|C{0.95cm}|C{0.95cm}|C{0.95cm}|C{0.95cm}|C{0.95cm}|C{0.95cm}|C{0.95cm}|C{0.95cm}|C{0.95cm}|}
\hline
\textbf{1} & \textbf{2} & \textbf{3} & \textbf{4} & \textbf{5} & \textbf{6} & \textbf{7} & \textbf{8} & \textbf{9} & \textbf{10} \\ \hline
\textbf{D} & \textbf{B} & \textbf{C} & \textbf{A} & \textbf{B} & \textbf{D} & \textbf{A} & \textbf{C} & \textbf{D} & \textbf{A} \\ \hline
\hline
\textbf{11} & \textbf{12} & \textbf{13} & \textbf{14} & \textbf{15} & \textbf{16} & \textbf{17} & \textbf{18} & \textbf{19} & \textbf{20} \\ \hline
\textbf{B} & \textbf{C} & \textbf{A} & \textbf{D} & \textbf{B} & \textbf{C} & \textbf{A} & \textbf{D} & \textbf{B} & \textbf{C} \\ \hline
\end{tabular}
\end{center}

\vspace{1pt}
\noindent\begin{minipage}[t]{0.58\linewidth}
\textbf{PART II. True/False Questions}\\[1.5pt]
\renewcommand{\arraystretch}{1.15}
\begin{tabular}{|c|c|c|c|c|}
\hline
\textbf{Q} & \textbf{a)} & \textbf{b)} & \textbf{c)} & \textbf{d)} \\ \hline
\textbf{1} & \textbf{T} & \textbf{T} & \textbf{F} & \textbf{T} \\ \hline
\textbf{2} & \textbf{F} & \textbf{F} & \textbf{T} & \textbf{T} \\ \hline
\textbf{3} & \textbf{T} & \textbf{F} & \textbf{T} & \textbf{F} \\ \hline
\textbf{4} & \textbf{F} & \textbf{T} & \textbf{T} & \textbf{F} \\ \hline
\end{tabular}
\end{minipage}\hfill
\begin{minipage}[t]{0.40\linewidth}
\textbf{PART III. Short-Answer Questions}\\[1.5pt]
\renewcommand{\arraystretch}{1.15}
\begin{tabular}{|c|c|c|c|}
\hline
\textbf{Q} & \textbf{Answer} & \textbf{Q} & \textbf{Answer} \\ \hline
\textbf{1} & $7$ & \textbf{4} & $0.28$ \\ \hline
\textbf{2} & $5$ & \textbf{5} & $6$ \\ \hline
\textbf{3} & $9$ & \textbf{6} & $900$ \\ \hline
\end{tabular}
\end{minipage}

\vspace{5pt}
\noindent\textbf{PART IV. Free-Response Summary:}\\[2pt]
- \textbf{Problem 1:} a) $\vec{BK} = -\dfrac{3}{4}\vec{AB} + \dfrac{1}{4}\vec{AC}$; $\vec{BD} = -\vec{AB} + \dfrac{1}{3}\vec{AC}$. b) $\vec{BK} = \dfrac{3}{4}\vec{BD} \Rightarrow B, K, D$ are collinear.\\
- \textbf{Problem 2:} a) $M(2, 0)$, minimum value is $8$. b) Circumcenter $I(1, 2)$.\\
- \textbf{Problem 3:} a) $\vec{F} = \vec{F}_1 + \vec{F}_2$, $|\vec{F}| = 7000\text{ N}$. b) $\vec{F}_1 = (5000, 0)$, $\vec{F}_2 = (-4000, 4000\sqrt{3})$, $\vec{F} = (1000, 4000\sqrt{3})$, $\alpha \approx 81.8^\circ$.

\vspace{2pt}
\begin{center}
  {\large\bfseries\color{ink} DETAILED SOLUTIONS}
\end{center}
\vspace{0pt}

\begin{multicols}{2}
\raggedcolumns
\small

\textbf{\color{accent}PART I. MULTIPLE-CHOICE}

\textbf{Question 1.} The nonzero vectors collinear with $\vec{AB}$ formed by the vertices of parallelogram $ABCD$ are $\vec{AB}, \vec{BA}, \vec{CD}, \vec{DC}$ ($4$ vectors). Choice \textbf{D}.

\textbf{Question 2.} In regular hexagon $ABCDEF$ with center $O$, segment $AB$ is parallel, identically oriented, and equal in length to $OC$, so $\vec{AB} = \vec{OC}$. Choice \textbf{B}.

\textbf{Question 3.} $\vec{MN} + \vec{PQ} = (\vec{MQ} + \vec{QN}) + (\vec{PN} + \vec{NQ}) = \vec{MQ} + \vec{PN} + (\vec{QN} + \vec{NQ}) = \vec{MQ} + \vec{PN}$. Choice \textbf{C}.

\textbf{Question 4.} $\vec{u} = \vec{AB} - \vec{AC} = \vec{CB} \Rightarrow |\vec{u}| = BC = \sqrt{3^2 + 4^2} = 5$. Choice \textbf{A}.

\textbf{Question 5.} $\vec{v} = \vec{BA} + \vec{BC} = \vec{BD}$. Since $ABCD$ is a rhombus of side $a$ with $\widehat{ABC} = 60^\circ$, $\triangle ABC$ is equilateral ($AC = a$), and diagonal $BD = 2 \cdot \dfrac{a\sqrt{3}}{2} = a\sqrt{3}$. Choice \textbf{B}.

\textbf{Question 6.} For centroid $G$ of $\triangle ABC$ and any point $O$: $\vec{OA} + \vec{OB} + \vec{OC} = 3\vec{OG} + (\vec{GA} + \vec{GB} + \vec{GC}) = 3\vec{OG}$. Choice \textbf{D}.

\textbf{Question 7.} $M$ lies on segment $AB$ with $MA = 3MB \Rightarrow AM = \dfrac{3}{4}AB$ and $\vec{AM}$ has the same direction as $\vec{AB}$, so $\vec{AM} = \dfrac{3}{4}\vec{AB}$. Choice \textbf{A}.

\textbf{Question 8.} $\vec{AM} = \vec{AB} + \vec{BM} = \vec{AB} + \dfrac{1}{2}\vec{BC} = \vec{AB} + \dfrac{1}{2}\vec{AD}$. Choice \textbf{C}.

\textbf{Question 9.} $\vec{OA} = (2, -1)$, $\vec{u} = (-3, 4) \Rightarrow \vec{v} = (2 - 3, -1 + 4) = (-1, 3)$. Choice \textbf{D}.

\textbf{Question 10.} $\vec{c} = 2(1, -2) - (-3, 5) = (2 + 3, -4 - 5) = (5, -9)$. Choice \textbf{A}.

\textbf{Question 11.} $x_G = \dfrac{1 - 3 + 5}{3} = 1$, $y_G = \dfrac{4 + 2 - 3}{3} = 1 \Rightarrow G(1, 1)$. Choice \textbf{B}.

\textbf{Question 12.} $\vec{AB} = (3, -6)$, $\vec{AC} = (m + 1, 2)$. Collinear $\Leftrightarrow 3(2) - (-6)(m + 1) = 0 \Leftrightarrow 6m + 12 = 0 \Leftrightarrow m = -2$. Choice \textbf{C}.

\textbf{Question 13.} $\vec{u} \parallel \vec{v} \Leftrightarrow 6m - (-4)(-3) = 0 \Leftrightarrow 6m = 12 \Leftrightarrow m = 2$. Choice \textbf{A}.

\textbf{Question 14.} Since $\triangle ABC$ is right isosceles at $A$, $\widehat{B} = 45^\circ$. Thus $(\vec{AB}, \vec{BC}) = 180^\circ - 45^\circ = 135^\circ$. Choice \textbf{D}.

\textbf{Question 15.} $\vec{AB} \cdot \vec{AC} = a \cdot a \cdot \cos 60^\circ = \dfrac{a^2}{2}$. Choice \textbf{B}.

\textbf{Question 16.} $\vec{a} + \vec{b} = (5, 3) \Rightarrow \vec{a} \cdot (\vec{a} + \vec{b}) = 2(5) + (-1)(3) = 7$. Choice \textbf{C}.

\textbf{Question 17.} $\vec{AB} = (-6, 8) \Rightarrow AB = \sqrt{(-6)^2 + 8^2} = 10$. Choice \textbf{A}.

\textbf{Question 18.} $\vec{u} \perp \vec{v} \Leftrightarrow 2(m + 1) + 3(-4) = 0 \Leftrightarrow 2m - 10 = 0 \Leftrightarrow m = 5$. Choice \textbf{D}.

\textbf{Question 19.} $\cos(\vec{x}, \vec{y}) = \dfrac{1\sqrt{3} + \sqrt{3}\cdot 1}{\sqrt{1+3}\sqrt{3+1}} = \dfrac{2\sqrt{3}}{4} = \dfrac{\sqrt{3}}{2} \Rightarrow (\vec{x}, \vec{y}) = 30^\circ$. Choice \textbf{B}.

\textbf{Question 20.} $W = F \cdot s \cdot \cos 30^\circ = 80 \cdot 20 \cdot \dfrac{\sqrt{3}}{2} = 800\sqrt{3}\text{ J}$. Choice \textbf{C}.

\columnbreak

\textbf{\color{accent}PART II. TRUE / FALSE}

\textbf{Question 1.}
a) Median $AM = \dfrac{a\sqrt{3}}{2} \Rightarrow |\vec{AG}| = \dfrac{2}{3}AM = \dfrac{a\sqrt{3}}{3} \Rightarrow$ \textbf{True}.\\
b) $\vec{AB} + \vec{AC} = 2\vec{AM} = 2 \cdot \dfrac{3}{2}\vec{AG} = 3\vec{AG} \Rightarrow$ \textbf{True}.\\
c) $(\vec{GA}, \vec{GB}) = 120^\circ$, $GA = GB = \dfrac{a\sqrt{3}}{3} \Rightarrow \vec{GA}\cdot\vec{GB} = \dfrac{a^2}{3}\cos 120^\circ = -\dfrac{a^2}{6} \ne \dfrac{a^2}{6} \Rightarrow$ \textbf{False}.\\
d) $|\vec{GB} - \vec{GC}| = |\vec{CB}| = a \Rightarrow$ \textbf{True}.

\textbf{Question 2.}
a) Since $M$ is the midpoint of $BC$, $\vec{AM} = \dfrac{1}{2}\vec{AB} + \dfrac{1}{2}\vec{AC} \Rightarrow$ \textbf{False}.\\
b) Since $N$ lies on $AC$ with $3NA = 2NC$, $AN = \dfrac{2}{5}AC \Rightarrow \vec{AN} = \dfrac{2}{5}\vec{AC} \ne \dfrac{2}{3}\vec{AC} \Rightarrow$ \textbf{False}.\\
c) $\vec{BN} = \vec{AN} - \vec{AB} = -\vec{AB} + \dfrac{2}{5}\vec{AC} \Rightarrow$ \textbf{True}.\\
d) $\vec{MN} = \vec{AN} - \vec{AM} = \dfrac{2}{5}\vec{AC} - \left(\dfrac{1}{2}\vec{AB} + \dfrac{1}{2}\vec{AC}\right) = -\dfrac{1}{2}\vec{AB} - \dfrac{1}{10}\vec{AC} \Rightarrow$ \textbf{True}.

\textbf{Question 3.}
a) $B$ is the midpoint of $AD \Rightarrow D = 2B - A = (7, 10) \Rightarrow$ \textbf{True}.\\
b) $\vec{AB} = (3, 4)$, $\vec{BC} = (5, 0)$. Although $BA = BC = 5$, $\vec{AB} \cdot \vec{BC} = 15 \ne 0$, so $\angle B \ne 90^\circ \Rightarrow$ \textbf{False}.\\
c) $\vec{AC} = (8, 4) \Rightarrow AC = \sqrt{64 + 16} = \sqrt{80} = 4\sqrt{5} \Rightarrow$ \textbf{True}.\\
d) $\vec{AB} \cdot \vec{AC} = 3(8) + 4(4) = 40 \ne 20 \Rightarrow$ \textbf{False}.

\textbf{Question 4.} Place axes at $A(0, 0), B(6, 0), C(6, 8), D(0, 8)$.\\
a) $\vec{AC} = (6, 8)$, $\vec{BD} = (-6, 8) \Rightarrow \vec{AC} \cdot \vec{BD} = -36 + 64 = 28 \ne 0 \Rightarrow$ \textbf{False}.\\
b) $\vec{AB} \cdot \vec{AC} = 6(6) + 0(8) = 36 \Rightarrow$ \textbf{True}.\\
c) $\vec{AC} + \vec{BD} = (0, 16) = 2\vec{AD} \Rightarrow |\vec{AC} + \vec{BD}| = 16 \Rightarrow$ \textbf{True}.\\
d) $\vec{AB} + \vec{AC} = (12, 8) \Rightarrow |\vec{AB} + \vec{AC}| = \sqrt{144 + 64} = 4\sqrt{13} \ne 14 \Rightarrow$ \textbf{False}.

\textbf{\color{accent}PART III. SHORT-ANSWER}

\textbf{Question 1.} $|\vec{u}| = |\vec{CB}| = BC = \sqrt{5^2 + 8^2 - 2\cdot 5\cdot 8\cos 60^\circ} = \sqrt{49} = \mathbf{7}$.

\textbf{Question 2.} $\vec{AI} = \dfrac{1}{2}\vec{AM} = \dfrac{1}{4}(\vec{AB} + \vec{AC})$. Hence $\vec{BI} = \vec{AI} - \vec{AB} = -\dfrac{3}{4}\vec{AB} + \dfrac{1}{4}\vec{AC}$. Thus $x = -\dfrac{3}{4}, y = \dfrac{1}{4} \Rightarrow P = 8\left(\dfrac{1}{4}\right) - 4\left(-\dfrac{3}{4}\right) = 2 + 3 = \mathbf{5}$.

\textbf{Question 3.} $\vec{AD} = \vec{BC} = (2, -4) \Rightarrow D = (1, -1)$. Thus $T = 4(1) - 5(-1) = \mathbf{9}$.

\textbf{Question 4.} $\vec{AB} = (3, 4)$, $\vec{AC} = (-3, 4)$. $\cos\widehat{BAC} = \dfrac{-9 + 16}{5 \cdot 5} = \dfrac{7}{25} = \mathbf{0.28}$.

\textbf{Question 5.} $\vec{u} = (2 - k, 2 + 3k)$. $\vec{u} \perp \vec{c} \Leftrightarrow 5(2 - k) + 1(2 + 3k) = 0 \Leftrightarrow 12 - 2k = 0 \Rightarrow k = \mathbf{6}$.

\textbf{Question 6.} $W = F \cdot s \cdot \cos 30^\circ = 40\sqrt{3} \cdot 15 \cdot \dfrac{\sqrt{3}}{2} = \mathbf{900}\text{ J}$.

\textbf{\color{accent}PART IV. FREE-RESPONSE}

\textbf{Problem 1.}
a) Since $I$ is the midpoint of $BC \Rightarrow \vec{AI} = \dfrac{1}{2}(\vec{AB} + \vec{AC})$.\\
Since $K$ is the midpoint of $AI \Rightarrow \vec{AK} = \dfrac{1}{2}\vec{AI} = \dfrac{1}{4}\vec{AB} + \dfrac{1}{4}\vec{AC}$.\\
Hence $\vec{BK} = \vec{AK} - \vec{AB} = -\dfrac{3}{4}\vec{AB} + \dfrac{1}{4}\vec{AC}$.\\
Since $AD = \dfrac{1}{3}AC \Rightarrow \vec{AD} = \dfrac{1}{3}\vec{AC}$, so $\vec{BD} = \vec{AD} - \vec{AB} = -\vec{AB} + \dfrac{1}{3}\vec{AC}$.\\
b) We see $\vec{BK} = \dfrac{3}{4}\left(-\vec{AB} + \dfrac{1}{3}\vec{AC}\right) = \dfrac{3}{4}\vec{BD}$.\\
Since $\vec{BK}$ and $\vec{BD}$ are parallel, points $B, K, D$ are collinear.

\textbf{Problem 2.}
a) Let $M(x, 0) \in Ox$. Then $\vec{MA} = (-2 - x, 6)$ and $\vec{MB} = (6 - x, 2)$.\\
Thus $\vec{MA} + \vec{MB} = (4 - 2x, 8) \Rightarrow |\vec{MA} + \vec{MB}| = \sqrt{(4 - 2x)^2 + 64} \ge 8$.\\
Equality holds when $4 - 2x = 0 \Leftrightarrow x = 2$. Thus $M(\mathbf{2}, \mathbf{0})$ and the minimum value is $\mathbf{8}$.\\
b) Let the circumcenter be $I(a, b)$. Then $IA^2 = IB^2 = IC^2$:\\
$\begin{cases} (a + 2)^2 + (b - 6)^2 = (a - 6)^2 + (b - 2)^2 \\ (a - 6)^2 + (b - 2)^2 = (a + 3)^2 + (b + 1)^2 \end{cases} \Leftrightarrow \begin{cases} 16a - 8b = 0 \\ -18a - 6b = -30 \end{cases}$\\
$\Leftrightarrow \begin{cases} b = 2a \\ 3a + b = 5 \end{cases} \Leftrightarrow \begin{cases} a = 1 \\ b = 2 \end{cases}$. Hence the circumcenter is $I(\mathbf{1}, \mathbf{2})$.

\textbf{Problem 3.}
a) The resultant force acting on the ship is $\vec{F} = \vec{F}_1 + \vec{F}_2$.\\
Its magnitude is $|\vec{F}| = \sqrt{F_1^2 + F_2^2 + 2F_1 F_2 \cos 120^\circ}$\\
$= \sqrt{5000^2 + 8000^2 + 2\cdot 5000\cdot 8000\cdot \left(-\dfrac{1}{2}\right)} = 1000\sqrt{49} = \mathbf{7000}\text{ N}$.\\
b) With $Ox$ along $\vec{F}_1$, we have $\vec{F}_1 = (\mathbf{5000}, \mathbf{0})$.\\
Since $(\vec{F}_1, \vec{F}_2) = 120^\circ$, $\vec{F}_2 = (8000\cos 120^\circ, 8000\sin 120^\circ) = (\mathbf{-4000}, \mathbf{4000\sqrt{3}})$.\\
Therefore $\vec{F} = \vec{F}_1 + \vec{F}_2 = (\mathbf{1000}, \mathbf{4000\sqrt{3}})$.\\
The angle $\alpha$ between $\vec{F}$ and $\vec{F}_1$ satisfies $\cos\alpha = \dfrac{\vec{F} \cdot \vec{F}_1}{|\vec{F}||\vec{F}_1|} = \dfrac{1000 \cdot 5000}{7000 \cdot 5000} = \dfrac{1}{7} \Rightarrow \alpha \approx \mathbf{81.8}^\circ$.

\end{multicols}
\end{document}
